

\lussec 1. Introduction

In view of its renormalization properties 
[\cite{WilsonFlow},\cite{RenFlow}], and since
its application in lattice gauge theory is technically straightforward,
the Yang--Mills gradient flow allows the dynamics of non-Abelian
gauge theories to be probed in many interesting ways.
The flow can be used for accurate scale setting, for example,
and it provides an understanding of how exactly
the topological (instanton) sectors emerge in
the continuum limit of lattice QCD
[\cite{WilsonFlow}]\kern1pt\sfn{$\dagger$}{%
In lattice gauge theory, where the right-hand side of the flow
equation coincides with the gradient of the Wilson plaquette action,
the flow is often referred to as the ``Wilson flow''. 
The term ``gradient flow'' is
used here both in the continuum theory and on the lattice.}.
Moreover, observables at positive flow time 
are natural quantities to consider for non-perturbative
renormalization and step scaling 
[\cite{StepScaling}--\cite{FritzschRamos}].

Matter fields may or may not be included in the flow equations.
In this paper, a fairly trivial extension of the flow to the 
quark fields in QCD is considered,
where the flow equation for the
gauge field is unchanged, while the evolution of 
the quark fields as a function of the flow time is determined
by a gauge-covariant heat equation (see sect.~2).
The theoretical analysis and practical implementation of the flow 
is not significantly complicated by
the inclusion of the quark fields. In particular,
following ref.~[\cite{RenFlow}], the renormalization of
correlation functions of gauge-invariant local fields at
positive flow times can be shown to require no more than
a multiplicative renormalization of
the time-dependent quark fields once
the parameters of QCD are renormalized as usual.

For illustration, two applications of the extended flow 
are worked out in this paper,
one being a new strategy for the
calculation of the axial-current renormalization constant
in lattice QCD
and the other a computation of the chiral condensate
essentially
through the evaluation of the expectation
value of the scalar quark density at positive flow time.
In both cases, the use of the flow proves to be
technically attractive.
The chiral condensate, for example, is easily obtained with
high precision, because no additive renormalization is required.

In the next two sections, the extension of the gradient
flow to the quark fields is discussed in the continuum theory.
Since the flow equations respect chiral symmetry,
the correlation functions of the time-dependent fields
satisfy simple chiral Ward identities.
As explained in sect.~4, these allow
the correlation functions to be related to the 
physics of the light pseudo-scalar mesons.
Then follows a more technical part of the paper
(sects.~5--7), where the flow is set up in the framework of 
lattice QCD. In particular, the associated field theory
in 4+1 dimensions is shown to admit a well-defined
local lattice regularization. 
The viability of the proposed applications of the 
flow in numerical lattice QCD is finally demonstrated 
through a sample calculation in 2+1 flavour QCD (sects.~8,9).

